\documentclass[a4paper]{report}
% Shared preamble for the English probability and statistics documents.
% Based on the original course preamble. Compile with pdfLaTeX.
\usepackage[top=1.5cm,bottom=2cm,left=2cm,right=2cm]{geometry}
\usepackage[T1]{fontenc}
\usepackage{lmodern}
\usepackage[utf8]{inputenc}
\usepackage[english]{babel}

\usepackage{amsmath}
\usepackage{amssymb}
\usepackage{amsthm}
\usepackage{mathrsfs}
\usepackage{mathtools}
\usepackage{etoolbox}
\usepackage{tikz}
\usetikzlibrary{shapes.multipart,arrows.meta}
\usepackage{xcolor}
\usepackage{enumerate}
\usepackage[shortlabels]{enumitem}
\usepackage{hyperref}
\usepackage{version}
\usepackage{bbold}
\usepackage{fontawesome5}
\usepackage{fancyhdr}
\lfoot{}
\usepackage{natbib}
\usepackage{graphicx}
\usepackage{pdfpages}
\usepackage{authblk}
    
\usepackage{blindtext}
\usepackage{multicol}
\usepackage{collcell}
\usepackage{wrapfig}
\usepackage{blkarray}
\usepackage{cancel}
\usepackage{appendix}
\usepackage{float}
\makeatletter
\input{tex/assets/slashbox.sty}
\makeatother
\hypersetup{unicode,colorlinks,citecolor=red,filecolor=black,
  linkcolor=red,urlcolor=cyan}
\setlength{\columnsep}{1cm}
\pagenumbering{gobble}

\def\numtext#1{%
  \ifcase#1\or one\or two\or three\or four\or five\or six\or seven\or eight\or nine\or ten\or
  eleven\or twelve\or thirteen\or fourteen\or fifteen\or sixteen\or seventeen\or eighteen\or nineteen\or twenty\or Lots\fi}

\newcommand{\Byte}[1]{%
  \def\nodecontents{}%
  \foreach[count=\l] \k in {#1} {\xappto{\nodecontents}{\noexpand\nodepart{\numtext{\l}}\k}}%
  \node[rectangle split,rectangle split parts=8,rectangle split horizontal,draw] at (2,2){\nodecontents};}

\newcommand{\D}{{\mathrm{d}}}
\newcommand{\cE}{{\mathcal{E}}}
\newcommand{\cG}{{\mathcal{G}}}
\newcommand{\cX}{{\mathcal{X}}}
\newcommand{\cY}{{\mathcal{Y}}}
\newcommand{\cB}{{\mathcal{B}}}
\newcommand{\cF}{{\mathcal{F}}}
\newcommand{\cO}{{\mathcal{O}}}
\newcommand{\cS}{{\mathcal{S}}}
\newcommand{\cT}{{\mathcal{T}}}
\newcommand{\given}{\, \big| \,}
\newcommand{\cL}{{\mathscr{L}}}
\newcommand{\cD}{{\mathcal{D}}}
\newcommand{\cA}{{\mathscr{A}}}
\newcommand{\cP}{{\mathscr{P}}}
\newcommand{\cR}{{\mathscr{R}}}
\newcommand{\A}{{\mathcal A}}
\newcommand{\one}{\mathbb{1}}
\newcommand{\mathbbm}[1]{\mathbb{#1}}

\providecommand{\C}{}
\renewcommand{\C}{\mathbb C}
\newcommand{\R}{\mathbb R}
\newcommand{\Q}{\mathbb Q}
\newcommand{\N}{\mathbb N}
\newcommand{\T}{\mathbb T}
\newcommand{\F}{\mathbb F}
\renewcommand{\P}{\mathbb P}
\newcommand{\Z}{\mathbb Z}
\newcommand{\E}{\mathbb {E}}
\newcommand{\Pbb}[1]{\P\left( #1\right)}
\newcommand{\Ebb}[1]{\E\left[ #1\right]}

\newcommand{\lb}{\left (}
\newcommand{\rb}{\right )}
\newcommand{\lbb}{\left [}
\newcommand{\rbb}{\right ]}
\newcommand{\labs}{\left |}
\newcommand{\rabs}{\right |}
\newcommand{\lbrb}[1]{\lb #1 \rb}
\newcommand{\lbbrbb}[1]{\lbb#1\rbb}
\newcommand{\labsrabs}[1]{\labs#1\rabs}
\newcommand{\lbbrb}[1]{\lbb#1\rb}
\newcommand{\lbrbb}[1]{\lb#1\rbb}
\newcommand{\langlerangle}[1]{\langle#1\rangle}
\newcommand{\lbcurly}{\{}
\newcommand{\rbcurly}{\}}
\newcommand{\curly}[1]{\lbcurly#1\rbcurly}
\newcommand{\abs}{\labsrabs}
\newcommand{\ind}[1]{1_{#1}}
\newcommand{\indep}{\perp \!\!\! \perp}
\renewcommand{\d}{\, \mathrm{d}}

\newtheorem{theorem}{Theorem}
\newtheorem{problemCnt}{Lemma}
\theoremstyle{definition}
\newtheorem{problem}[problemCnt]{Problem}
\newtheorem{problembook}[problemCnt]{\faBook \,  Problem}
\newtheorem{problempencil}[problemCnt]{\faPencil* \,  Problem}
\newtheorem{problemstar}[problemCnt]{\faStar\,  Problem}
\newtheorem{problemstarpencil}[problemCnt]{\faStar \, \faPencil* \,   Problem}
\newtheorem{probleminteresting}[problemCnt]{\faSmile[regular] \,   Problem}
\newtheorem{solution}{Solution}

\newcommand{\fwcell}[1]{\makebox[\arraystretch\normalbaselineskip]{$#1$}}
\newenvironment{enumeratei}{\begin{enumerate}[(i)]}{\end{enumerate}}
\newenvironment{enumeratea}{\begin{enumerate}[(a)]}{\end{enumerate}}
\newenvironment{enumeraten}{\begin{enumerate}[\upshape 1.]}{\end{enumerate}}

% Compatibility helpers needed by a few files in the combined archive.
\AtBeginDocument{\noindent{\footnotesize English translation prepared with ChatGPT (OpenAI, GPT-6), 26 October 2026, for translation of the Bulgarian text into English and LaTeX formatting.
The Bulgarian originals remain the reference; translation may contain errors.}\par\smallskip}
\newcommand{\MissingGraphic}[1]{%
  \begin{center}\fbox{\parbox{0.82\linewidth}{\centering Missing source image: \texttt{#1}}}\end{center}}



\setlength{\emergencystretch}{3em}
\allowdisplaybreaks
\begin{document}


\noindent Probability and Statistics, 21.01.2024 \hfill Test 2

\hrulefill

\renewcommand{\P}{\mathbb{P}}
\renewcommand{\D}{\mathbb{D}}

\begin{it}
The exact grading formula depends on the results. As an approximation, use 2 plus the points earned. Time allowed: 3 hours. Good luck.

$m,n,k$ denote nonnegative integers.



\end{it}

\begin{problem}
    Let $X$ have distribution function
\[
F_X(x) = \begin{cases}
    0, & \text{for }x\leq0;\\
    ax^2+bx, & \text{for }0<x\leq1;\\
    1, & \text{for }x>1.
\end{cases}
\]
    \begin{enumerate}
        \item (0.25 points) Is $X$ a continuous random variable? If so, derive its density.
        \item (0.25 points) Find $a,b$ if additionally $\P(X>0.75)=0.5$.
        \item (0.25 points) If the expectation of a random variable $Y$ exists, then
        \[\mathbb{E}[Y]=\int_0^\infty\mathbb{P}(Y>t)\D t-\int_{-\infty}^0\mathbb{P}(Y\leq t)\D t.\]
        Prove this identity when $Y$ is continuous with density $f_Y$.
        \item (0.25 points) Use the formula to find $\E X$.
    \end{enumerate}
\end{problem}

\begin{solution}
    \begin{enumerate}
        \item Yes: it is continuous because it takes uncountably many values. It is also absolutely continuous because the given distribution function is differentiable:
            \[f_X(x) = \frac{d}{dx}F_X(x) = \begin{cases}
            0, & \text{for }x\leq0;\\
            2ax+b, & \text{for }0<x<1;\\
            0, & \text{for }x>1.
            \end{cases} = (2ax + b)\mathbb{1}_{\{0<x<1\}}\]

        \item We are also given $\P(X>0.75)=1-F_X(0.75)=0.5$. From the definition, $1=\P(X<1)=F_X(1)=a+b$. Solve the system:
        \[
        \begin{array}{|l} 1-F_X\left(\frac{3}{4}\right) = \frac{1}{2} \\ F_X(1) = 1\end{array} \Leftrightarrow
        \begin{array}{|l} F_X\left(\frac{3}{4}\right) = \frac{1}{2} \\ F_X(1) = 1\end{array} \Leftrightarrow
        \begin{array}{|l} \frac{9}{16}a + \frac{3}{4}b = \frac{1}{2} \\ a+b = 1\end{array}\]
        Thus $(a,b)=\left(\frac16,\frac56\right)$.

        \item The alternative expression for the expectation follows by interchanging the order of integration. We show that the right side equals the left. Consider the first summand:
        \[\int_0^\infty\P(Y>t)dt = \int_0^\infty\left(\int_t^\infty f_Y(y)dy\right)dt = \int\limits_{\{0 < y < t\}}f_Y(y)d(t, y) = \int_0^\infty\left(\int_0^yf_Y(y)dt\right)dy =\]
        \[\int_0^\infty\left(\int_0^ydt\right)f_Y(y)dy = \int_0^\infty yf_Y(y)dy\]

        Similarly, for the second summand,
        \[-\int_{-\infty}^0\P(Y\le t)dt = \int_{-\infty}^0yf_Y(y)dy\]

        completing the proof.

        Formally, we used Fubini's theorem, which allows the order of integration to be interchanged when the integral over the region is absolutely convergent. Explain why this applies here.
        \item
        \begin{align*}
            \E[X]
            &= \int_0^\infty\P(X>t)dt - \int_{-\infty}^0\P(X < t)dt \\
            &= \int_0^\infty(1-F_X(t))dt - \underbrace{\int_{-\infty}^0F_X(t)dt}_{=0} \\
            &= \int_0^1(1-ax^2-bx)dx + \underbrace{\int_1^\infty(1-F_X(x))dx}_{=0} \\
            &= 1 - \frac{a}{3} - \frac{b}{2} \\
            &= \frac{19}{36}
        \end{align*}

        To check the arithmetic, compare the result with the usual calculation $\int_0^1xf_X(x)\,dx$.

    \end{enumerate}
\end{solution}

\begin{problem}
    Let $U_1,U_2,U_3\sim U(0,1)$ be independent, and let $U_{(1)},U_{(2)},U_{(3)}$
    be the values $U_1,U_2$
    and $U_3$ in increasing order.

\textit{For example, if $U_1=0.23$, $U_2=0.88$, $U_3=0.1$, then
$U_{(1)}=0.1,$ $ U_{(2)}=0.23,$ $ U_{(3)}=0.88$. }
    \begin{enumerate}
        \item (0.25 points) Find the densities and expectations of $U_{(1)}:=\min\{U_1,U_2,U_3\}$ and $U_{(3)}:=
        \max\{U_1, U_2, U_3\}$.
        \item (0.25 points) Is it true that
        $U_{(2)}\sim U(0,1)$? If yes, prove it;
        otherwise find its distribution.
        \item (0.25 points) Find the variance
        of $X:=-\ln(U_1)$.
    \end{enumerate}
    Let $X_1,X_2,X_3,X_4\sim Exp(1)$
    be independent.
    \begin{enumerate}
         \setcounter{enumi}{3}
        \item (0.25 points) Find the densities,
        expectations and variances of
        $S_2:=X_1+X_2$, $S_3:=X_1+X_2+X_3$ and $S_4:=X_1+\dots+X_4$.
        \item (0.5 points) Let $Y_k:={S_k}/{
        S_4}$. Prove that $Y_{(k)}$ and $U_{(k)
}$ have the same distribution for $k=1,2,3$.

    \end{enumerate}
\end{problem}

\begin{solution}
    \begin{enumerate}
        \def\U#1{U_{(#1)}}
        \def\F#1{F_{U_{(#1)}}}
        \def\f#1{f_{U_{(#1)}}}
        \def\Ex#1{E[U_{(#1)}]}
        \item Proceed in the standard way (for $0<t<1$; outside this interval the result is immediate):
        \begin{align*}
            \F1(t) &= \P(\U1<t) = 1-\P(\U1>t) = 1 - \P(U_1>t)^3 = 1-(1-t)^3 \\
            \Rightarrow \f1(t) &= \frac{d}{dt}\F1(t) = -3(1-t)^2\cdot(-1) = 3(1-t)^2 \\
            \Rightarrow \Ex1 &= \int_0^1t\cdot3(1-t)^2dt = 3\int_0^1(t^3 - 2t^2 + t)dt = \dots = \frac{1}{4}
        \end{align*}
        \begin{align*}
            \F3(t) &= \P(\U3<t) = \P(U_1 < t)^3 = t^3 \\
            \Rightarrow \f3(t) &= 3t^2 \\
            \Rightarrow \Ex3 &= \int_0^1t\cdot3t^2dt = 3\int_0^1t^3dt =\dots=\frac{3}{4}
        \end{align*}

        \item It is worth considering why the initial guess is misleading, but directly:
        \begin{align*}
            \F2(t) &= \P(\U2 < t) \\
                &=\P\left(\text{at least two of the three independent values are }<t\right)\\
                &= \P(Bin(3, t) \ge 2) \\
                &= \binom{3}{2}t^2(1-t) + t^3 \\
                &= 3t^2-2t^3 \\
            \Rightarrow \f2(t) &= 6t - 6t^2 = 6t(1-t)
        \end{align*}

        \item Write $U:=U_1$. We have $X:=-\ln(U)=:g(U)\in[0,+\infty)$. The function $g$ is strictly decreasing and hence invertible; its inverse $h(x):=g^{-1}(x)=e^{-x}$ is differentiable. Changing variables,

        \[f_X(x)=\underbrace{f_U(h(x))}_{=1}\cdot|\underbrace{h'(x)}_{-e^{-x}}|=e^{-x}\text{ for }0<x.\]

        Hence $X\sim Exp(1)$, and we may use $\Ebb{X}
        =DX=1$ or calculate directly:
        \begin{align*}
            \E[X] &= \int_{0}^{\infty}xe^{-x}dx = \dots = 1 \\
            \E[X^2] &= \int_{0}^{\infty}x^2e^{-x}dx = \dots = 2 \\
            \Rightarrow \D[X] &= 2 - 1 = 1
        \end{align*}

        \item Recall that a sum of independent identically distributed exponential random variables is gamma-distributed. If $X_i\sim Exp(\lambda)$ are independent, then
        \[S_n := \sum_{i = 1}^nX_i \sim \Gamma(n, \lambda)\]

        Thus, in general and (on the right) in our particular case, for $t>0$:
        \begin{align*}
            f_{S_n}(t) &= \frac{\lambda^nt^{n-1}e^{-\lambda t}}{(n-1)!} &&= \frac{t^{n-1}}{(n-1)!}e^{-t} \\
            \E[S_n] &= \frac{n}{\lambda} &&= n \\
            \D[S_n] &= \frac{n}{\lambda^2} &&= n
        \end{align*}

        The expectation and variance follow directly from independence. Even without using the gamma distribution, the density can be derived from the convolution formula for a sum of two independent random variables:
        \begin{align*}
            f_{S_2}(t) &\stackrel{t>0}{=} \int_{-\infty}^{+\infty}\underbrace{f_{X_1}(x)}_{=e^{-x}\cdot\mathbb{1}_{\{0<x\}}}\underbrace{f_{X_2}(t-x)}_{e^{-(t-x)}\cdot\mathbb{1}_{\{0 < t - x\}}}dx = \int_0^te^{-x}e^{x-t}dx = e^{-t}\int_0^tdx &&= te^{-t}\\
            f_{S_3}(t) &\stackrel{t>0}{=} \int_{-\infty}^{+\infty}\underbrace{f_{S_2}(x)}_{=xe^{-x}\cdot\mathbb{1}_{\{0<x\}}}\underbrace{f_{X_3}(t-x)}_{e^{-(t-x)}\cdot\mathbb{1}_{\{0 < t - x\}}}dx = \int_0^txe^{-x}e^{x-t}dx = e^{-t}\int_0^txdx &&= \frac{t^2}{2}e^{-t}\\
            f_{S_4}(t)&\stackrel{t>0}{=}\dots\text{ by induction}.
        \end{align*}

        \item For the last part, write
        \[
        Y_k = \frac{X_1 + \dots +  X_k}{(X_1 + \dots + X_k) +
        (X_{k+1} + \dots + X_{4})} = \frac{S_k}{S_k + \widetilde{S}_{4-k}},
        \]
        where $\widetilde S_{4-k}$ is a copy of
        $S_{4-k}$ independent of $S_k$. Then use the standard method
        for the density of a transformation (a two-dimensional change of variables, its Jacobian, etc.).
    \end{enumerate}
\end{solution}

\begin{problem}
\begin{enumerate}
    \item (0.25 points)
In the manufacture of specialized automotive parts, the lengths of certain components are measured,
and deviations from the recommended length are modelled with mean 0 and variance 9.
A component is acceptable if its deviation from the recommended length is less than a threshold \(s\).
What percentage of length errors fall within the acceptable range when \(s=1\)?
Find \(s\) such that 95\% of components meet the acceptance criterion.

\item (0.5 points) Let $X\sim N(1,1)$. To accuracy 0.01,
find $\E|X|$.
\end{enumerate}
\end{problem}

\begin{solution}
    \begin{enumerate}
        \item
        \item \begin{align*}
            \E[|X|] &= \int_{-\infty}^{+\infty}|x|f_X(x)dx \stackrel{t := x - 1}{=} \int_{-\infty}^{+\infty}|t+1|\underbrace{f_X(t+1)}_{f_Z(t)}dt \\
            &= \int_{-\infty}^{-1}-(t+1)f_Z(t)dt + \int_{-1}^{+\infty}(t+1)f_Z(t)dt \\
            &= \underbrace{-\int_{-\infty}^{-1}f_Z(t)dt}_{=-\Phi(-1)} + \underbrace{\int_{-1}^{+\infty}f_Z(t)dt}_{=1 - \Phi(-1)} - \underbrace{\int_{-\infty}^{-1}tf_Z(t)dt}_{t \mapsto -t} + \underbrace{\int_{-1}^{+\infty}tf_Z(t)dt}_{=\int_{-1}^{-1} + \int_1^\infty} \\
            &=1-2\underbrace{\Phi(-1)}_{1-\Phi(1)}+2\int_1^\infty tf_Z(t)\,dt+\underbrace{\int_{-1}^1tf_Z(t)\,dt}_{=0\text{ (odd)}}\\
            &= 2\Phi(1) - 1 + 2J
        \end{align*}

        It remains to calculate:
        \begin{align*}
            J := \int_1^\infty tf_Z(t)dt &= \int_1^\infty t\cdot\frac{1}{\sqrt{2\pi}}e^{-\frac{t^2}{2}}dt = -\frac{1}{\sqrt{2\pi}}\int_1^\infty e^{-\frac{t^2}{2}}d\left(-\frac{t^2}{2}\right) = -\frac{1}{\sqrt{2\pi}}\left[e^{-\frac{t^2}{2}}\right]_{t=1}^{t\to+\infty} = \frac{1}{\sqrt{2\pi e}}
        \end{align*}

        Finally:
        \[\E[|X|] = -1 + 2\underbrace{\Phi(1)}_{\approx 0.841} + \underbrace{\frac{2}{\sqrt{2\pi e}}}_{\approx 0.484} \approx 1.166\]
    \end{enumerate}
\end{solution}

\begin{problem}
The plane is ruled by parallel lines spaced $2t$ apart. A needle of length $2l\leq2t$ is dropped
at a random position.
\begin{enumerate}
    \item (0.5 points) What is the probability that the needle touches a line?

\textit{Hint: Parametrize by the position of the needle's centre and its angle to the lines.}

    \item (0.25 points) Suppose $l=t$.
    What minimum number of independently dropped needles makes the probability that at least half cross a line at least 95\%?
Explain how and why this procedure can estimate $\pi$.

\item \textit{(0.25 points) Bonus: Solve part 1 when $l>t$.}
\end{enumerate}
\end{problem}

\begin{solution}
    This is known as Buffon's needle problem.
    \begin{enumerate}
        \item Let
        \begin{align*}
            X&:=\text{distance from the needle's centre to the nearest line}&&\sim U(0,t)\\
            \theta&:=\text{acute angle between the needle and the nearest line}&&\sim U\left(0,\frac\pi2\right).
        \end{align*}

        The needle crosses the line if and only if $X\leq l\sin\theta$, so we seek:
        \[p := \P(X \le l\sin\theta) = \int_0^{\frac{\pi}{2}}\int_{0}^{l\sin\theta}\frac{2}{\pi t}dtd\theta = \frac{2l}{\pi t}\]
        See the linked \href{https://en.wikipedia.org/wiki/Buffon\%27s_needle_problem}{solution of Buffon's needle problem}.
        \item Since $l=t$, $p=2/\pi$. Model the needle drops as Bernoulli trials: $S_n:=\sum_{i=1}^nX_i$, where the independent $X_i\sim Ber(p)$ indicate success when a needle crosses a line. We seek an estimate of $n$ such that
        \[\P\left(S_n \ge \frac{1}{2}n\right) \ge 95\%\]

        Use the CLT. First note $\mu:=\E[X_1]=p=2/\pi$ and $\sigma:=\sqrt{\D[X]}=\sqrt{p(1-p)}$. We require:

        \[\P\left(\frac{S_n - n\mu}{\sigma\sqrt{n}} \ge \frac{\frac{n}{2} - n\mu}{\sigma\sqrt{n}}\right)
        =\P\left(\frac{S_n-n\mu}{\sigma\sqrt n}\geq\frac{\sqrt n(\frac12-\mu)}{\sigma}\right)\stackrel{\text{CLT}}{\approx}
        \P\left(Z \ge \underbrace{\frac{\sqrt{n}\left(\frac{1}{2} - \mu\right)}{\sigma}}_{=: C_n}\right) \ge 95\%\]

        \begin{align*}
            1 - \Phi\left(C_n\right) &\ge 95\% \\
            \Phi(C_n)&\leq5\%&&\Phi\text{ is increasing, as is }\Phi^{-1}\\
            C_n &\le \Phi^{-1}(5\%)\\
            \frac{\sqrt n(\frac12-\mu)}{\sigma}&\leq\Phi^{-1}(5\%)&&C_n<0\text{ and decreases with }n,\text{ since }\tfrac12<\mu\\
            \sqrt n&\geq\frac{\sigma\Phi^{-1}(5\%)}{\frac12-\mu}&&\text{note also }\Phi^{-1}(5\%)<0\\
            n &\ge \left(\frac{\sigma \Phi^{-1}(5\%)}{\frac{1}{2}-\mu}\right)^2
        \end{align*}

        Approximating:
        \begin{align*}
            \Phi^{-1}(5\%) &\approx -1.645 \\
            \sigma &\approx 0.481 \\
            \mu &\approx 0.637 \\
            \Rightarrow \left(\frac{\sigma \Phi^{-1}(5\%)}{\frac{1}{2}-\mu}\right)^2 &\approx 33.5387
        \end{align*}

        Finally, $n\geq34$.

        The (strong) law of large numbers explains the estimate of $\pi$. Since $\frac1nS_n\stackrel{\text{a.s.}}{\longrightarrow}\mu=2/\pi$, estimate $\pi$ by $2/(\frac1nS_n)$. Consider how large $n$ should be for a desired accuracy.

        \item See the long-needle case in the linked solution of \href{https://en.wikipedia.org/wiki/Buffon\%27s_needle_problem}{Buffon's needle problem}.
    \end{enumerate}
\end{solution}




























































\end{document}
